\documentclass[12pt]{article}
%ACENTOS
\usepackage[utf8]{inputenc}

%PAQUETES
\usepackage{amsfonts,amsmath,amssymb}
\usepackage{amsmath,latexsym}
\usepackage{graphics,color, xcolor, soulutf8, enumitem, float}
\usepackage{graphicx}
\usepackage[most]{tcolorbox}
\usepackage{hyperref} %Para colocar direcciones Web
\usepackage{indentfirst}
\usepackage{wrapfig}  %Figuras flotantes
\usepackage{subfig}   %Un ambiente Figura con varias imagenes
\usepackage{caption}  %Caption a subfig
\usepackage{multicol}

%MARGEN
\setlength\topmargin{-0.5in}
\addtolength\oddsidemargin{-0.7in}
\setlength{\textheight}{23cm}
\setlength{\textwidth}{17cm}

%VARIABLES
\newcommand{\R}{\mathbb R}

%COLORES
\newcommand{\red}[1]{\textcolor[rgb]{1.00,0.00,0.00}{#1}}
\newcommand{\blue}[1]{\textcolor[rgb]{0.00,0.00,1.00}{#1}}
\newcommand{\verde}[1]{\textcolor[rgb]{0.00,0.50,0.25}{#1}}

%OTROS
\def\refname{Referencias}%

\renewcommand{\baselinestretch}{1}
\renewcommand\familydefault{\sfdefault}

\begin{document}
{\footnotesize
	\begin{tabular}{cc}
		\hspace*{-1cm}
		\begin{minipage}[H]{10cm}
			\noindent UNIVERSIDAD CENTRAL DE VENEZUELA\\
			FACULTAD DE CIENCIAS\\
			ESCUELA DE COMPUTACIÓN
		\end{minipage}
		&
		\begin{minipage}[H]{10cm}
			CÁLCULO CIENTÍFICO (6109)\\
			SEMESTRE II-2025\\
			MARZO DE 2026
		\end{minipage} \\
\end{tabular}}
\vspace*{2mm}
\begin{center}
	{\Large \textbf{Proyecto: Curvas de Nivel}} \\
	\vspace*{3mm}
	{\normalsize Maximiliano Barboza \hspace{1.5cm} Yunior Moreno} \\
	{\footnotesize CI: 28.314.688 \hspace{2.15cm} CI: 26.465.131}
\end{center}

%----------------------------------------------------------
\section*{Motivación}
%----------------------------------------------------------

En múltiples áreas de la ciencia e ingeniería resulta necesario representar información tridimensional sobre un plano bidimensional. Una superficie definida por $f(x,y)=z$ con $f:\R^2\to\R$ se puede caracterizar parcialmente a través de sus \textit{curvas de nivel}: al interceptar la superficie con planos horizontales $z=k$ se obtienen curvas de contorno que, proyectadas sobre el plano $xy$ conforman dichas curvas de nivel \cite{penney94}.

En cartografía, por ejemplo, las curvas de nivel sirven para representar relieves topográficos. Cuando están muy juntas indican pendientes pronunciadas, y cuando están separadas el terreno es más plano. Pero más allá de visualizar el relieve, el \textbf{área entre curvas de nivel} también tiene uso práctico: en obras de construcción permite estimar volúmenes de excavación o relleno, y eso se traduce directamente en costos.

Este informe se centra en ese problema: \textit{calcular el área entre curvas de nivel}. Primero se trabaja con el caso de funciones, que es más sencillo, y después se pasa a curvas cerradas parametrizadas.

%----------------------------------------------------------
\section*{Área entre dos curvas (funciones)}
%----------------------------------------------------------

Si $p_n$ es el polinomio interpolante construido con los nodos $(x_i, f(x_i))$ $i=0,1,\ldots,n$ para una cierta función $f$ continua en $[a,b] = [x_0,x_n]$ tiene sentido plantear la aproximación
\[
\int_a^b f(x)\,dx \approx \int_a^b p_n(x)\,dx = I_n(f)
\]
donde $I_n(f)$ suele denominarse \textit{cuadratura numérica}. La idea es sencilla: como integrar $f$ directamente puede ser difícil, se reemplaza por un polinomio cuya integral sí se puede calcular sin problemas \cite{cheney08}.

Si se dispone de dos funciones $f$ y $g$ definidas en $[a,b]$ el área de la región comprendida entre sus gráficas viene dada por
\[
\text{Area} = \int_a^b |f(x) - g(x)|\,dx
\]
En la práctica las curvas pueden tener cambios bruscos de pendiente o cruzarse en varias zonas del intervalo, así que conviene dividir $[a,b]$ en subintervalos y construir un interpolante independiente en cada uno. De esa forma se controla mejor la aproximación y se evitan oscilaciones innecesarias. En la implementación decidimos conservar esta idea y permitir divisiones de tramo marcadas por el usuario, para que una variación local no obligue a ajustar toda la curva con un único interpolante.

%----------------------------------------------------------
\subsection*{Requerimientos teóricos}
%----------------------------------------------------------

\begin{enumerate}
	%==== ITEM 1 ====
	\item \textbf{Preferencia de base polinomial para integración.}

	El polinomio interpolante $p_n(x)$ se puede expresar en distintas bases del espacio $\Pi_n$: la base canónica $\{1,x,x^2,\ldots,x^n\}$ la base de Lagrange $\{L_0(x),L_1(x),\ldots,L_n(x)\}$ o la base de diferencias divididas de Newton. Como el polinomio interpolante es único sin importar la base, las tres representaciones dan el mismo valor $I_n(f)$ Pero a la hora de calcular integrales, la \textbf{base de Lagrange} resulta más conveniente \cite{burden02}.

	Al escribir
	\[
	p_n(x) = \sum_{i=0}^{n} f(x_i)\,L_i(x)
	\]
	la integral se convierte directamente en
	\[
	I_n(f) = \int_a^b p_n(x)\,dx = \sum_{i=0}^{n} f(x_i) \underbrace{\int_a^b L_i(x)\,dx}_{A_i} = \sum_{i=0}^{n} A_i\,f(x_i)
	\]
	Lo importante aquí es que los coeficientes $A_i$ solo dependen de los nodos $x_0,x_1,\ldots,x_n$ y del intervalo $[a,b]$ pero no de la función $f$ Eso quiere decir que se calculan una sola vez y sirven para cualquier función. Con la base canónica o la de Newton esa separación no sale tan directa, porque los coeficientes mezclan los valores $f(x_i)$ de forma más complicada.

	%==== ITEM 2 ====
	\item \textbf{Método de los coeficientes indeterminados.}

	Hay otra forma de obtener los pesos de cuadratura sin necesidad de construir el polinomio interpolante: el \textbf{método de los coeficientes indeterminados} \cite{chapra07}. Se parte de que la fórmula tiene la forma
	\[
	I_n(f) = \sum_{i=0}^{n} A_i\,f(x_i)
	\]
	donde los $A_i$ son incógnitas. Para hallarlos, se exige que la fórmula sea \textit{exacta} para $f(x) = 1,\, x,\, x^2, \ldots, x^n$ lo cual genera un sistema de $n+1$ ecuaciones con $n+1$ incógnitas:
	\[
	\sum_{i=0}^{n} A_i\,x_i^k = \int_a^b x^k\,dx, \quad k=0,1,\ldots,n
	\]
	Este sistema siempre tiene solución única porque la matriz asociada es de tipo Vandermonde (invertible cuando los nodos son distintos). Lo bueno de este enfoque es que se obtienen los pesos resolviendo un sistema lineal, sin tener que construir el polinomio ni calcular las integrales de los $L_i(x)$ Además, es fácil verificar si la fórmula tiene un \textit{grado de exactitud} mayor al esperado: solo hay que probar si también es exacta para $f(x) = x^{n+1}, x^{n+2},\ldots$

	%==== ITEM 3 ====
	\item \textbf{Cuadratura con nodos equidistantes en $[0,1]$}

	Consideremos $x_0=0$, $x_1=\tfrac{1}{2}$ y $x_2=1$, con $y_i = f(x_i)$ para $i=0,1,2$ Buscamos $A_0, A_1, A_2$ tales que
	\[
	I_2(f) = A_0\,y_0 + A_1\,y_1 + A_2\,y_2
	\]
	Aplicando el método de coeficientes indeterminados, imponemos exactitud para $f(x) = 1,\, x,\, x^2$:
	\begin{align}
		f(x)=1:   &\quad A_0 + A_1 + A_2 = \int_0^1 1\,dx = 1 \label{eq:s1}\\[4pt]
		f(x)=x:   &\quad A_0\cdot 0 + A_1\cdot\tfrac{1}{2} + A_2\cdot 1 = \int_0^1 x\,dx = \tfrac{1}{2} \label{eq:s2}\\[4pt]
		f(x)=x^2: &\quad A_0\cdot 0 + A_1\cdot\tfrac{1}{4} + A_2\cdot 1 = \int_0^1 x^2\,dx = \tfrac{1}{3} \label{eq:s3}
	\end{align}

	Restando (\ref{eq:s3}) de (\ref{eq:s2}):
	\[
	\tfrac{1}{2}A_1 - \tfrac{1}{4}A_1 = \tfrac{1}{2} - \tfrac{1}{3} \quad\Longrightarrow\quad \tfrac{1}{4}A_1 = \tfrac{1}{6} \quad\Longrightarrow\quad A_1 = \frac{2}{3}
	\]
	Sustituyendo $A_1$ en (\ref{eq:s2}): $\tfrac{1}{2}\cdot\tfrac{2}{3} + A_2 = \tfrac{1}{2}$ de donde $A_2 = \tfrac{1}{2} - \tfrac{1}{3} = \tfrac{1}{6}$

	De (\ref{eq:s1}): $A_0 = 1 - \tfrac{2}{3} - \tfrac{1}{6} = \tfrac{1}{6}$

	Se obtiene entonces:
	\[
	\boxed{I_2(f) = \frac{1}{6}\,y_0 + \frac{2}{3}\,y_1 + \frac{1}{6}\,y_2}
	\]

	Los valores $(A_0,A_1,A_2)^T = \left(\tfrac{1}{6},\,\tfrac{2}{3},\,\tfrac{1}{6}\right)^T$ son los \textbf{pesos de cuadratura}, también conocidos como coeficientes de Newton-Cotes \cite{mathworld_newtoncotes}. Cada $A_i$ indica cuánto ``pesa'' el valor $f(x_i)$ en la aproximación de la integral. Esta fórmula resulta ser la conocida \textit{regla de Simpson} aplicada a $[0,1]$ \cite{cheney08, burden02}.

	\begin{itemize}
		\item \textbf{Exactitud para $f \in \Pi_3$} La fórmula $I_2(f)$ fue diseñada para ser exacta en $\Pi_2$ pues se usaron las condiciones $f=1,\,x,\,x^2$ Sin embargo, también resulta exacta para polinomios de grado 3. Como la integral y la cuadratura son operadores lineales, basta verificar para $f(x) = x^3$ (ya que $\{1,x,x^2,x^3\}$ es base de $\Pi_3$):
		\[
		I_2(x^3) = \frac{1}{6}\cdot 0^3 + \frac{2}{3}\cdot\left(\frac{1}{2}\right)^3 + \frac{1}{6}\cdot 1^3 = \frac{2}{3}\cdot\frac{1}{8} + \frac{1}{6} = \frac{1}{12} + \frac{2}{12} = \frac{1}{4}
		\]
		Por otro lado:
		\[
		\int_0^1 x^3\,dx = \frac{x^4}{4}\bigg|_0^1 = \frac{1}{4}
		\]
		Dado que $I_2(x^3) = \int_0^1 x^3\,dx$ y ya se tiene exactitud para $1,\,x,\,x^2$ se concluye por linealidad que $I_2(f) = \int_0^1 f(x)\,dx$ para todo $f\in\Pi_3$ $\blacksquare$

		Ese grado de exactitud extra viene de la simetría de los nodos respecto al centro del intervalo, y es algo particular de la regla de Simpson \cite{mathworld_simpson}.

		Para llevar esta fórmula a un intervalo general $[a,b]$ se hace el cambio $x = a + (b-a)\,u$ con $u\in[0,1]$ lo que da
		\[
		\int_a^b f(x)\,dx \approx \frac{b-a}{6}\left[f(a) + 4\,f\!\left(\frac{a+b}{2}\right) + f(b)\right]
		\]
		que es la regla de Simpson ``simple'' sobre $[a,b]$ Cuando se requiere más precisión, se divide $[a,b]$ en $N$ subintervalos (con $N$ par) y se aplica la regla en cada par consecutivo, dando origen a la \textbf{regla de Simpson compuesta} que se emplea en la implementación.
	\end{itemize}

	%==== ITEM 4 ====
	\item \textbf{Cuadratura con nodos no equidistantes.}

	Ahora tomemos $x_0=0$, $x_1=\tfrac{1}{3}$ y $x_2=1$. Aplicando nuevamente el método:
	\begin{align}
		f(x)=1:   &\quad A_0 + A_1 + A_2 = 1 \label{eq:n1}\\[4pt]
		f(x)=x:   &\quad \tfrac{1}{3}A_1 + A_2 = \tfrac{1}{2} \label{eq:n2}\\[4pt]
		f(x)=x^2: &\quad \tfrac{1}{9}A_1 + A_2 = \tfrac{1}{3} \label{eq:n3}
	\end{align}

	Restando (\ref{eq:n3}) de (\ref{eq:n2}):
	\[
	\left(\tfrac{1}{3} - \tfrac{1}{9}\right)A_1 = \tfrac{1}{2} - \tfrac{1}{3} \quad\Longrightarrow\quad \tfrac{2}{9}\,A_1 = \tfrac{1}{6} \quad\Longrightarrow\quad A_1 = \frac{3}{4}
	\]
	De (\ref{eq:n2}): $A_2 = \tfrac{1}{2} - \tfrac{1}{3}\cdot\tfrac{3}{4} = \tfrac{1}{2} - \tfrac{1}{4} = \tfrac{1}{4}$

	De (\ref{eq:n1}): $A_0 = 1 - \tfrac{3}{4} - \tfrac{1}{4} = 0$

	La nueva cuadratura queda:
	\[
	I_2(f) = \frac{3}{4}\,y_1 + \frac{1}{4}\,y_2
	\]

	Llama la atención que $A_0=0$ es decir, el valor $f(0)$ no aporta nada a la aproximación. Eso ya sugiere que estos nodos no están bien repartidos.

	Verifiquemos para $f(x) = x^3$:
	\[
	I_2(x^3) = \frac{3}{4}\cdot\left(\frac{1}{3}\right)^3 + \frac{1}{4}\cdot 1^3 = \frac{3}{4}\cdot\frac{1}{27} + \frac{1}{4} = \frac{1}{36} + \frac{9}{36} = \frac{10}{36} = \frac{5}{18}
	\]
	Mientras que $\int_0^1 x^3\,dx = \tfrac{1}{4} = \tfrac{9}{36} \neq \tfrac{5}{18}$

	\textbf{Conclusión:} esta nueva $I_2(f)$ \textit{no} es exacta para $f\in\Pi_3$ Al romper la simetría de los nodos se pierde el grado de exactitud extra que tenía Simpson. Queda claro que cómo se distribuyan los nodos afecta mucho la calidad de la cuadratura \cite{chapra07}.
\end{enumerate}


%----------------------------------------------------------
\section*{Área entre dos curvas cerradas (relaciones)}
%----------------------------------------------------------

Como se mencionó en la motivación, una curva de nivel no es necesariamente el gráfico de una función; más bien es una figura cerrada en el plano. Para calcular áreas en este caso hay que recurrir al cálculo vectorial, en particular al \textbf{teorema de Green} y a la \textbf{integral de línea} \cite{stewart12}.

\subsection*{Requerimientos teóricos}

\begin{enumerate}
	%==== ITEM 1 ====
	\item \textbf{Área de una figura cerrada a partir de mediciones discretas.}

	Sea $C$ una curva cerrada simple (sin autointersecciones) recorrida en sentido antihorario, y sea $D$ la región que encierra. El teorema de Green establece que \cite{stewart12, mathworld_green}
	\[
	\oint_C \bigl(P\,dx + Q\,dy\bigr) = \iint_D \left(\frac{\partial Q}{\partial x} - \frac{\partial P}{\partial y}\right) dA
	\]
	Si elegimos $P(x,y) = -y$ y $Q(x,y) = x$ entonces $\dfrac{\partial Q}{\partial x} - \dfrac{\partial P}{\partial y} = 1 + 1 = 2$ y la doble integral se reduce a $2\,\text{Area}(D)$ De aquí se obtiene la conocida fórmula
	\begin{equation}\label{eq:area_green}
	\text{Area}(D) = \frac{1}{2}\oint_C \bigl(x\,dy - y\,dx\bigr)
	\end{equation}

	Ahora bien, en la práctica no se dispone de la curva continua sino de un conjunto finito de mediciones $(x_i, y_i)$, $i=0,1,\ldots,n$, con la condición $(x_0,y_0) = (x_n,y_n)$, que garantiza el cierre de la figura. El procedimiento más directo para aproximar el área consiste en unir los puntos consecutivos mediante segmentos rectos (interpolación lineal a trozos). Sobre cada segmento que va de $(x_i, y_i)$ a $(x_{i+1}, y_{i+1})$ la contribución a la integral de línea es
	\[
	\frac{1}{2}\bigl(x_i\,y_{i+1} - x_{i+1}\,y_i\bigr)
	\]
	Sumando sobre todos los segmentos se llega a la \textbf{fórmula del cordón} (también llamada \textit{Shoelace formula}) \cite{mathworld_shoelace}:
	\begin{equation}\label{eq:shoelace}
	\text{Area} = \frac{1}{2}\left|\sum_{i=0}^{n-1}\bigl(x_i\,y_{i+1} - x_{i+1}\,y_i\bigr)\right|
	\end{equation}
	El valor absoluto se pone para que el área salga positiva sin importar el sentido de recorrido. Si la curva es un polígono, la fórmula da el resultado exacto; para curvas suaves, la precisión depende de qué tan bien los segmentos rectos se ajusten a la curva real.

	%==== ITEM 2 ====
	\item \textbf{Rol de la integral de línea y la parametrización.}

	La ecuación (\ref{eq:area_green}) expresa el área como una \textit{integral de línea} sobre la curva $C$ Para evaluarla, es necesario describir la curva mediante una parametrización: si $C$ está dada por las funciones $x(t)$ e $y(t)$ con $t\in[t_0, t_n]$ donde $(x(t_0),y(t_0)) = (x(t_n),y(t_n))$ entonces $dx = x'(t)\,dt$ y $dy = y'(t)\,dt$ y la integral de línea se convierte en una integral ordinaria \cite{stewart12}:
	\begin{equation}\label{eq:area_param}
	\text{Area}(D) = \frac{1}{2}\int_{t_0}^{t_n} \bigl(x(t)\,y'(t) - y(t)\,x'(t)\bigr)\,dt
	\end{equation}

	Esta transformación convierte un problema en el plano en uno de una sola variable. En nuestro caso tenemos puntos discretos $(x_i, y_i)$ sobre la curva, y a cada uno le asignamos un valor del parámetro $t_0 < t_1 < \cdots < t_n$ Lo más sencillo es tomar $t_i = i$ para $i=0,1,\ldots,n$ aunque también se podría usar la longitud de arco acumulada \cite{cheney08}.

	Con los pares $(t_i, x_i)$ y $(t_i, y_i)$ se construyen por separado dos interpolantes (o \textit{splines}): uno para $x(t)$ y otro para $y(t)$ De esta forma, lo que era una curva en el plano se convierte en dos problemas de interpolación en una variable, que es justo lo que ya se sabe hacer con las técnicas de la sección anterior.

	La evaluación de (\ref{eq:area_param}) requiere conocer las derivadas $x'(t)$ e $y'(t)$ Si se emplea el polinomio interpolante de Lagrange, estas derivadas se pueden obtener de manera analítica. Recordemos que
	\[
	x(t) = \sum_{i=0}^{n} x_i\,L_i(t), \qquad L_i(t) = \prod_{\substack{j=0\\j\neq i}}^{n} \frac{t - t_j}{t_i - t_j}
	\]
	Derivando respecto a $t$ y aplicando la regla del producto sobre cada $L_i$:
	\begin{equation}\label{eq:deriv_lagrange}
	x'(t) = \sum_{i=0}^{n} x_i\,L_i'(t), \qquad L_i'(t) = \sum_{\substack{k=0\\k\neq i}}^{n} \frac{1}{t_i - t_k}\prod_{\substack{j=0\\j\neq i\\j\neq k}}^{n} \frac{t - t_j}{t_i - t_j}
	\end{equation}
	La expresión es larga pero se puede programar sin mayor dificultad. Lo bueno es que no hace falta recurrir a diferenciación numérica (cocientes de diferencias finitas), que tiende a amplificar errores de redondeo \cite{burden02}.

	%==== ITEM 3 ====
	\item \textbf{Integración numérica en este contexto.}

	Una vez obtenidas las funciones interpolantes $x(t)$ e $y(t)$ (ya sea mediante polinomios de Lagrange por trozos o \textit{splines}) se requiere calcular la integral (\ref{eq:area_param}). Para ello se puede emplear cualquier regla de cuadratura; en particular, la \textbf{regla de Simpson compuesta} ofrece un buen balance entre sencillez y precisión \cite{burden02}. El integrando es
	\[
	h(t) = x(t)\,y'(t) - y(t)\,x'(t)
	\]
	donde las derivadas $x'(t)$ e $y'(t)$ se obtienen mediante la expresión (\ref{eq:deriv_lagrange})

	Un aspecto delicado es la elección entre un \textit{único} polinomio interpolante global de grado $n$ o una interpolación \textit{por trozos} con polinomios de grado bajo. Con pocos puntos de control ($n \leq 8$ aproximadamente), el polinomio global suele comportarse bien. Sin embargo, cuando el número de puntos crece, aparecen oscilaciones espurias entre los nodos (el conocido \textbf{fenómeno de Runge} \cite{cheney08}) que distorsionan tanto la curva interpolada como el cálculo del área.

	Para evitar ese problema se usa \textbf{interpolación local por ventanas}: en cada segmento paramétrico $[t_i, t_{i+1}]$ se toma una ventana de $d+1$ nodos vecinos (con $d=3$ típicamente) centrada alrededor del segmento y se arma un polinomio de Lagrange \textit{local} solo con esos nodos. Con ese polinomio se evalúan $x(t)$ $y(t)$ y se calculan las derivadas $x'(t)$ $y'(t)$ dentro del segmento y luego se aplica Simpson sobre $[t_i, t_{i+1}]$ Se suman las contribuciones de todos los segmentos y eso da el área total. Al mantener los polinomios en grado bajo se evitan las oscilaciones de Runge, y se puede trabajar con cualquier cantidad de puntos sin que la interpolación se vuelva inestable. En la implementación decidimos usar este enfoque local en lugar de un interpolante global, porque nos interesaba privilegiar estabilidad cuando el usuario coloca muchos nodos.

	Para calcular el \textbf{área entre dos curvas cerradas} $C_1$ y $C_2$ (por ejemplo, dos curvas de nivel consecutivas donde $C_2$ encierra a $C_1$), se calcula el área de cada región por separado y luego se obtiene la diferencia:
	\[
	\text{Area entre curvas} = \bigl|\text{Area}(D_2) - \text{Area}(D_1)\bigr|
	\]
	donde $D_1$ y $D_2$ son las regiones encerradas por $C_1$ y $C_2$ respectivamente. Es la misma idea de la Parte 1 donde se integraba $|f(x) - g(x)|$ solo que aquí cada ``función'' es en realidad el área que encierra una curva parametrizada.

	Además, la fórmula del cordón (\ref{eq:shoelace}) es un caso particular de todo esto: equivale a usar interpolación lineal ($d=1$) entre puntos consecutivos y calcular la integral de línea de forma exacta sobre cada segmento recto. Si en cambio se usa interpolación de grado mayor ($d=3$ por ejemplo), las curvas entre nodos ya no son rectas sino arcos suaves, y eso mejora la aproximación del área para figuras con contornos curvos, aunque requiere más cálculos.
\end{enumerate}


%----------------------------------------------------------
\section*{Descripción de los métodos numéricos empleados}
%----------------------------------------------------------

En la implementación computacional se utilizan los siguientes métodos numéricos:

\begin{itemize}
	\item \textbf{Interpolación mediante spline cúbico natural (área entre funciones).} Las curvas $f(x)$ y $g(x)$ se interpolan a partir de los nodos que el usuario coloca sobre la imagen, usando un spline cúbico con condición de frontera natural. Si el usuario define líneas de división, se construye un spline independiente en cada subintervalo con los nodos que le caen. Un spline cúbico es básicamente una función polinómica a trozos de grado~3 que pasa por todos los nodos y cuya primera y segunda derivada son continuas en los puntos de unión; que sea ``natural'' quiere decir que $S''(x_0) = S''(x_n) = 0$ en los extremos, así la curva no se fuerza a curvarse en los bordes \cite{burden02, cheney08}. En el código se usa la clase \texttt{CubicSpline} de \texttt{scipy.interpolate} \cite{scipy_cubicspline} con \texttt{bc\_type="natural"}; por dentro, la rutina arma y resuelve un sistema tridiagonal para obtener los coeficientes de cada trozo cúbico, lo cual es eficiente ($O(n)$) por la estructura de banda de la matriz \cite{burden02}.

	\item \textbf{Interpolación local de Lagrange por ventanas (área entre curvas cerradas).} Para las curvas cerradas parametrizadas se interpola con Lagrange de forma local: en cada segmento $[t_i, t_{i+1}]$ se toma una ventana de $d+1$ nodos vecinos (con $d=3$) y se arma un polinomio solo con esos nodos. Así se obtienen $x(t)$ $y(t)$ y sus derivadas $x'(t)$ $y'(t)$ de forma analítica con (\ref{eq:deriv_lagrange}) sin necesidad de un polinomio global de grado alto que podría oscilar \cite{cheney08}. La aplicación también ofrece la fórmula del cordón (\ref{eq:shoelace}) como método alternativo, útil para comparar resultados o cuando los puntos están tan cerca que la aproximación poligonal ya es suficiente. Probamos dejar este método alternativo en la aplicación, y observamos que sirve como contraste rápido frente al efecto de suavizar la curva.

	\item \textbf{Regla de Simpson compuesta.} Las integrales se aproximan con Simpson compuesta usando $N = 200$ subintervalos por defecto. Como se vio en la parte teórica, esta regla es exacta para polinomios de grado $\leq 3$ y tiene un error de orden $O(h^4)$ donde $h$ es el ancho de cada subintervalo \cite{mathworld_simpson, burden02}. La fórmula que se implementó es:
	\[
	\int_a^b f(x)\,dx \approx \frac{h}{3}\left[f(x_0) + f(x_N) + 4\sum_{\substack{i=1 \\ i\text{ impar}}}^{N-1} f(x_i) + 2\sum_{\substack{i=2 \\ i\text{ par}}}^{N-2} f(x_i)\right]
	\]
\end{itemize}


\bibliography{biblio}
\bibliographystyle{plain}

\end{document}
